\section{Conclusion}
\label{sec:conclusion}

Caches for diffusion transformers decide when to skip the block stack and what to write instead, but the fillers they
use are open loop: the errors of skipped steps are integrated into the latent and never corrected. \pcgs closes this
loop at no cost in network evaluations. It treats each skipped gap as a smoothing problem, estimates the residuals of
the gap with the RTS smoother of the forecaster's own state-space model once the next step is computed, and
re-integrates the gap through the output layers. Combined with the \sea gate and a Kalman residual forecaster, it
improves fidelity to the 50-step sampler at every budget we tested on \flux, with gains over reuse that grow with the
budget. \TODO{One sentence on video.} Promising directions are to choose the computed steps jointly with the
corrector, rather than with a gate designed for reuse, and to apply the correction to caches that operate on features
inside the block stack.

\paragraph{Limitations.}
Matched \nfe is not matched latency: \pcgs adds output-layer evaluations (about 0.3\% of one transformer evaluation per
image) and 0.02\,GiB of memory over \kalsea, and the filter adds elementwise work on every step.
\TODO{Latency once measured.} Our metrics measure fidelity to the 50-step sampler, which is the goal of a cache but
not the same as sample quality. The filter's band statistics are accumulated with atomic GPU additions, so runs are not
bitwise reproducible (0.03\,dB between two runs of the same configuration). All design decisions were made on a
development set of COCO captions and evaluated on held-out DrawBench prompts, but our image results cover a single
model. At a computed step, \pcgs attaches the residual evaluated at the uncorrected latent to the corrected one, an
approximation that is accurate only for small corrections; at the smallest budget, with gaps of up to twelve steps, its
gain is modest (0.31\,dB at 9 \nfe).
